\documentclass[11pt]{article}
\usepackage{elicitation}
\begin{document}

\packettitle{Collator's Notes}{For the person running the round, and for nobody else}

\begin{warning}
\textbf{Do not send this document to a respondent.} It names what the form deliberately withholds,
including how many rows the book leaves empty and which they are. A respondent who reads it is no
longer independent and their form cannot be used.
\end{warning}

\section{What the form withholds, and why}

The respondent form does not say how many rows the book leaves empty, or which. An earlier draft
did: it stated that the book's chapters rest on the claim that \textit{five} Traditions govern
nothing any Step consumes. A respondent reads that before filling in the grid, and with only
twelve rows a count is most of the way to the answer. It was removed on 10 August 2026 for the
same reason the form declines to name the two decisive cells.

\begin{note}
\textbf{For your reference and not for theirs:} the book leaves \textbf{five} rows empty --- T4,
T6, T7, T9 and T10 --- and fills 35 of the 96 cells.
\end{note}

\section{How many respondents}

Three is the useful minimum. With three you can report, per cell, whether zero, one, two or three
marked it, and the empty-rows claim becomes a statement about agreement rather than about one
person. Two is worth having. One is worth having and should be described as one.

\section{Collecting the forms}

Store each completed form in the repository as
\texttt{research/governance-elicitation-<initials>.md}. The script reads every file matching that
pattern: a markdown table whose first column is the Tradition and whose remaining eight columns
are the resources in the order given on the form. Blank and \texttt{0} both mean empty.

Record in \texttt{research/SOURCES.md} who filled each one in, what they were told, and whether
they had seen the book. \textbf{Whether the respondent has read Part Four is the single most
important thing to record}, because a reader who has is not independent and their form should be
reported separately rather than pooled.

\section{Running the comparison}

\begin{quote}\ttfamily\small
cd <repository root>\\
python3 model/elicitation\_compare.py -{}-self-test\\
python3 model/elicitation\_compare.py
\end{quote}

Run the self-test before sending any forms out. It passed on 10 August 2026.

\section{What the script reports, and why in that order}

\begin{enumerate}[leftmargin=*,itemsep=5pt]
\item \textbf{Which rows each respondent left empty}, and whether they match the book's five. This
      is first because Chapters 16, 17 and 18 all rest on it. If a respondent leaves four or six,
      or a different five, the script says Chapter 18 is contradicted, in those words.
\item \textbf{Cell-by-cell agreement out of 96}, with the two kinds of disagreement separated:
      cells the respondent filled that the book leaves empty, and cells the book fills that the
      respondent leaves empty. Those are different mistakes and one agreement rate hides it.
\item \textbf{The consequences}, with Chapters 16 and 17 recomputed on the respondent's matrix.
\end{enumerate}

\begin{note}
\textbf{Read kappa, not the raw agreement rate.} The book's matrix is sparse, 35 of 96 cells, so a
respondent marking at random with the same density would agree about 55 per cent of the time by
luck. The script prints Cohen's kappa alongside and the raw rate should not be quoted alone.
\end{note}

\section{Bare marks versus strengths}

A respondent may give only zeros and ones. The rerun then substitutes the book's own magnitude in
every cell both agree is non-zero, and 0.5 in cells the respondent filled and the book did not.
The script states that substitution in its output every time. A bare-marks respondent is testing
the \textbf{sparsity pattern} rather than the magnitudes, which is the more important half, so
this is not a degraded result.

\section{The analysis is preregistered}

\texttt{model/elicitation\_compare.py} was written before any form came back. That is the point of
it, and it is why the script should not be edited once forms start arriving. If it turns out to
need a change, make the change, say so in \texttt{research/progress-log.md}, and report both the
original and the revised analysis.

\section{Afterwards}

Whatever the result, propagate it through every layer listed in the synchronisation rule in
\texttt{CLAUDE.md}: the manuscript, the appendix, the paper, the primer, the ledgers, the plans,
the README, the progress log, \texttt{HANDOFF.md}, \texttt{AGENT\_VERIFY.md}, and all three PDFs.
Then re-run the full sequence in \texttt{README.md}.

\begin{warning}
If the result contradicts Chapter 18, the chapter says it will say so. Hold it to that.
\end{warning}

\end{document}
